\chapter{Análisis de datos}\label{cap:analisis}

El análisis parte de lo dicho en la sección~\ref{sec:alcance} del capítulo~\ref{cap:intro} y define la media muestral
\begin{equation}\label{eq:media}
  \bar{x} = \frac{1}{n}\sum_{i=1}^{n} x_i .
\end{equation}

\section{Varianza}\label{sec:varianza}

La varianza se calcula sobre la media de la ecuación~\eqref{eq:media}:
\begin{equation}\label{eq:varianza}
  s^2 = \frac{1}{n-1}\sum_{i=1}^{n} \left(x_i - \bar{x}\right)^2 .
\end{equation}

Otro párrafo largo en una sola línea para que el capítulo tenga cuerpo suficiente: se menciona la ecuación~\eqref{eq:varianza}, se vuelve a la introducción en la sección~\ref{sec:alcance} y se anticipa la conclusión del capítulo~\ref{cap:conclusion}, de modo que hay referencias hacia atrás y hacia adelante en el mismo párrafo.
